\chapter{Avanza adecuadamente en el diseño}
\label{cap:avance}

\section{El backlog y su registro}
\label{sec:backlog}

El backlog se estableció en el Informe de Etapa 1, con veinticinco ítems en siete épicas, y
se amplió el 4 de septiembre de 2026 con nueve tareas de proceso (H1 a H9), pendientes de
adopción. Su definición está en el repositorio (\href{run:../evidence/research/backlog/README.md}{backlog}) y su
ejecución en la planilla compartida con el profesor guía (\href{https://docs.google.com/spreadsheets/d/1FXSB_2WqGL2TbOgw698gQEBTbL3nV7qXFypcflpYgEw/edit?gid=313395138}{planilla}).
Cada ítem lleva la categoría con que la pauta mide el avance. Dos cambios de
planificación del 4 de septiembre se declaran: E3 pasó de la Etapa 4 a la 2, porque es una
decisión de diseño con su insumo (D1) terminado, y H8 pasó de la Etapa 3 a la 2, porque es el
diseño experimental que G2 compromete. El cuadro \ref{tab:backlog} da el estado al 8 de
octubre de 2026. Planilla y backlog enlazados conservan estados anteriores: G2 y E3 sin empezar,
H8 pendiente y fecha antigua. Su sincronización sigue abierta (anexo
\ref{anexo:procedencia}).

{\scriptsize
\setlength{\tabcolsep}{3pt}
\begin{longtable}{L{0.55cm}L{6.8cm}L{1.45cm}L{1.2cm}L{2.2cm}}
\caption[Registro del backlog entero al 8 de octubre de 2026.]{Registro del backlog entero,
con el estado al 8 de octubre de 2026; en negrita, las tareas de diseño. Categorías: impl.,
implementación; doc., documentación; pub., publicación.}\label{tab:backlog}\\
\toprule
ID & Tarea & Categoría & Etapa & Estado \\
\midrule
\endfirsthead
\multicolumn{5}{l}{\scriptsize\emph{Cuadro \ref{tab:backlog}, continuación}}\\
\toprule
ID & Tarea & Categoría & Etapa & Estado \\
\midrule
\endhead
\bottomrule
\endfoot
A1 & Estimar posiciones en un panel de un periodo & impl. & E1 & listo \\
A2 & Trayectorias multiperiodo en un marco común & impl. & E1 & listo \\
A3 & Valores iniciales declarados y reproducibles & impl. & E1 & listo \\
A4 & Validar entrada y reportar faltas de semilla & impl. & E2 & sin empezar \\
A5 & Reconciliar marco de exportación y optimizador & impl. & E2 & sin empezar \\
B1 & Cuatro búsquedas por optimizadores modernos & impl. & E1 & listo \\
B2 & Evidencia controlada a ajustes igualados & prueba & E3 & sin empezar \\
C1 & Comparar cada rutina con el estado intermedio & prueba & E1 & listo \\
C2 & Puntuar el estado terminal con la referencia & prueba & E1 & listo \\
C3 & Comparar mapas con amplitud y reflexiones & prueba & E1 & listo \\
C4 & Regenerar el brazo dinámico estadounidense & prueba & E3 & sin empezar \\
C5 & Repetir la verificación en otra máquina & prueba & E3 & sin empezar \\
C6 & Acordar el umbral de fidelidad por panel & decisión & E2 & del propietario \\
D1 & Sensibilidad a periodo, semilla y compilación & prueba & E1 & listo \\
D2 & Registro de sensibilidad y su razón & doc. & E4 & sin empezar \\
E1 & Intervalos para las medianas de partido & impl. & E1 & listo \\
E2 & Elegibilidad, multiplicidad y robustez & impl. & E1 & listo \\
\textbf{E3} & \textbf{Diseño de periodización y discretización} & diseño & E2 & por ratificar \\
E4 & El caso como validación del instrumento & doc. & E4 & sin empezar \\
F1 & Compilar y ejecutar desde un clon limpio & pub. & E4 & sin empezar \\
F2 & Publicar el repositorio con licencia & pub. & E4 & sin empezar \\
G1 & Capítulos 1 a 3 y anexo de backlog & doc. & E1 & listo \\
\textbf{G2} & \textbf{Capítulo de diseño y diseño experimental} & diseño & E2 & hecho \\
G3 & Implementación y resultados preliminares & doc. & E3 & sin empezar \\
G4 & Documento completo y conclusiones & doc. & E4 & sin empezar \\
H1 & Descomposición de objetivos & decisión & E2 & propuesta \\
H2 & Acordar RNF01 antes de diseñar pruebas & decisión & E2 & del propietario \\
H3 & Puerta de paridad de entradas & impl. & E3 & por adoptar \\
H4 & Suite de pruebas como artefacto & impl. & E3 & por adoptar \\
H5 & Registro de números en la tesis & doc. & E3, E4 & por adoptar \\
H6 & Tabla de retroalimentación por etapa & doc. & E2 a E4 & hecho en E2 \\
H7 & Presupuesto de páginas & decisión & E2 & supuesto \\
\textbf{H8} & \textbf{Pruebas exhaustivas diseñadas antes de correr} & diseño & E2 & escrito; aprobar \\
H9 & Congelar dos días antes de cada etapa & proceso & E2 a E4 & al congelar \\
\end{longtable}
}

\section{Tareas de diseño frente al backlog entero}
\label{sec:tareas-diseno}

El registro base contiene treinta y cuatro ítems; siete propuestas I1--I7 de la Etapa 3
se agregaron localmente el 8 de octubre y se detallan en el anexo \ref{anexo:procedencia}.
Siguen pendientes de adopción, con responsables y horas propuestas. En el registro base, tres son de diseño (G2, E3, H8); los tres están
comprometidos en la Etapa 2 y tienen su artefacto escrito aquí (cuadro
\ref{tab:tareas-diseno}). Los tres diseños están disponibles; la ratificación de E3 y la
aprobación prevista en H8 siguen abiertas. No se confunde artefacto escrito con tarea aceptada. Los diez ítems de implementación se miden en las
etapas 3 y 4; seis están listos desde la Etapa 1.

\begin{table}[H]
\centering
\caption{Tareas de diseño de la Etapa 2: artefacto escrito y aceptación pendiente.}
\label{tab:tareas-diseno}
\footnotesize
\begin{tabular}{L{3.2cm}L{6.0cm}L{3.2cm}}
\toprule
Tarea & Artefacto & Estado \\
\midrule
G2, capítulo de diseño y diseño experimental & capítulo \ref{cap:diseno}: vistas, decisiones
con su razón, atributos de calidad, verificación por niveles, diseño experimental, mejora de
rendimiento, trazabilidad; anexos \ref{anexo:figuras} a \ref{anexo:rendimiento} & hecho \\
\addlinespace
H8, diseño de las pruebas exhaustivas & sección \ref{sec:experimental}; anexo
\ref{anexo:paneles} & escrito; aprobación y RNF01 del propietario \\
\addlinespace
E3, periodización y discretización & sección \ref{sec:algoritmos} y anexo \ref{anexo:chile}:
diseños de periodización y regla de discretización; D-16; RD-8 & escrito; ratificación pendiente \\
\bottomrule
\end{tabular}
\end{table}

I1--I6 añadirían seis ítems de implementación y llevarían la base de diez a dieciséis
solo después de su adopción; I7 es documentación. El cálculo $10/16$ daría 62,5\,\%; es
condicional: los seis ítems listos darían $6/16$ (37,5\,\%), y E3 exige ocho cierres.
La adopción revisará solapamientos con A4, A5 y H4. I6 conserva wmay como antecedente; se fija original
2004 como referente externo de la tesis (cuadro \ref{tab:protocolos}).

Las demás tareas de la Etapa 2 no son de diseño y su estado se declara sin atenuantes. A4 y
A5, de implementación en el motor fiel, no se iniciaron, porque en esta etapa no se
implementó código nuevo en los motores (sección \ref{sec:cambios}); su reprogramación se
decide con el profesor guía. C6 y H2 esperan la decisión del propietario; H1 tiene una
propuesta (sección \ref{sec:propuesta-chile}); H6 está hecha en el capítulo \ref{cap:retro};
H7 quedó resuelta como supuesto de trabajo, y H9 se ejecuta al congelar este informe.
